\section{Certified solver islands and theory combination}

\subsection{Export a fragment, not an untyped approximation of Lean}

A \emph{solver island} is a selected subproblem whose translation and proof
reconstruction are supported end to end. Boolean and fixed-width bit-vector
islands can use the existing SAT-oriented facilities. A suitable first-order
island can use Lean-SMT and its reconstruction path; a clause-based island can
use Duper \cite{lean-grind,lean-smt,duper}. The island may prove a lemma over
shared interface terms rather than the complete original goal.

Each export has a translation ledger: Lean expressions, their solver symbols,
their types and operations, and proof obligations connecting the encoding to
the original proposition. Universe instantiation, monomorphization, typeclass
resolution, and auxiliary assertions must have justified transformations.
The Lean-SMT work illustrates why that preprocessing is a substantial part of
an integration rather than mere serialization \cite{lean-smt}.

For an island under assumptions $A_1,\ldots,A_k$, reconstruct a proof of
$A_1\to\cdots\to A_k\to L$. Apply it to the appropriate scoped proofs to
insert $L$ into the main graph. An external \code{unsat} answer without a
supported proof is \code{unknown}. An external \code{sat} answer is a model
proposal, not necessarily a counterexample to the Lean goal. Unsupported proof
rules and unsupported theories must be detected and reported, not silently
trusted.

\subsection{Do not import unsound combination assumptions}

A generic textbook theory-combination procedure is not automatically valid
for every Lean type. Finite domains, dependent types, casts, and overlapping
arithmetic operators complicate disjointness and cardinality assumptions.
Forge avoids relying on a blanket combination theorem: plugins exchange
\emph{proved interface facts}, and every final step is checked in Lean.
This is sound even when incomplete.

When a plugin wants $u=v$ to reconcile a model, it may request a case split or
ask another plugin to prove it. It cannot simply merge those terms. Similarly,
an arithmetic atom that abbreviates $xy$ must carry a definition or equality
proof; treating it as an unrelated real variable gives only a relaxation, not
a bidirectional equivalence of problems.

\subsection{Two trust profiles}

The default target is \textbf{kernel-checked replay}: ordinary proof terms or a
proved certificate-checker theorem whose acceptance computation is justified
without trusting an external search program. The intended trusted assumptions
are Lean's ordinary accepted foundations and explicitly audited library
axioms. Replay should not introduce \code{sorryAx}, a newly assumed theorem,
or a hidden oracle axiom.

A separate, explicitly named \textbf{native-checking profile} may run a
certificate checker using native compilation for performance. This has a
different trust story. The reference warns that \code{native\_decide} adds the
Lean compiler to the trusted part, whereas \code{decide\_cbv} does not require
trust in the code generator \cite{lean-tactics}. Do not summarize both profiles
as ``only the kernel is trusted.'' A SAT certificate and a theorem about its
checker are not, by themselves, a complete description of how the checker was
actually evaluated.

The precise existing implementation of a delegated tactic must be audited
under the pinned version. In strict mode, disable a backend when its available
replay path requires the native profile. Certificate size, kernel-replay time,
and axiom dependencies belong in the result record alongside search time.

\section{Concrete Lean integration plan}

\subsection{Use the current extension interface}

The inspected API has \code{Lean.Meta.Grind.SolverExtension} and
\code{registerSolverExtension}. The latter registers an extension during
initialization and allocates its state. Its methods include internalization,
new equality and disequality notifications, an action, consistency-related
checks, and a model-combination hook. State access, term marking, fact insertion,
and proof explanations are also exposed \cite{grind-types}.

The API skeleton below records actual conceptual hook names; the handler
bodies and Forge state are \textbf{design pseudocode}, not a compiled module.

\begin{Code}
initialize polyExt := registerSolverExtension makePolyState
initialize polyExt.setMethods
  (internalize := collectSupportedArithmetic)
  (newEq       := updateCanonicalAtoms)
  (newDiseq    := wakeNonzeroGuards)
  (action      := boundedCertificateAction)
  (checkInv    := checkInternalBookkeeping)
\end{Code}

Make the adapter a narrow, version-specific module. It reifies supported
expressions with a valuation map, requests a certificate, reconstructs a
proof, and submits a proved fact through the appropriate insertion service.
The check of internal data-structure invariants is not the mathematical
soundness proof of a certificate. Name and test those two responsibilities
separately.

Structural induction is better implemented above the local saturation loop:
it changes the goal shape and creates dependent subgoals. The resulting branch
contexts can each host a \grind\ session and its plugin state. This two-level
architecture keeps the existing theory engine local while the outer graph
tracks recursor applications and generalized lemmas.

\subsection{Module boundaries and required contracts}

\begin{center}
\begin{tabularx}{\textwidth}{@{}p{0.27\textwidth}YY@{}}
\toprule
Proposed module & Responsibility & Acceptance contract\\
\midrule
\code{Forge.Context} & Scoped expression, fact, and obligation graph & No unproved facts; valid dependency scopes\\
\code{Forge.GrindAdapter} & Version-specific extension hooks & Restorable state; exact expression/proof transport\\
\code{Forge.Demand} & Retrieval, instantiation, tabling & Every application elaborates and has acyclic proof parents\\
\code{Forge.Poly.Reify} & Typed polynomial representation & Denotation equality for each reified term\\
\code{Forge.Poly.Cert} & Cone and box certificate checking & Proved checker soundness or explicit proof reconstruction\\
\code{Forge.Induct} & Recursor selection and generalization & Goal refinement through an actual induction theorem\\
\code{Forge.Synth} & Invariants and existential candidates & Base/step or witness specification discharged\\
\code{Forge.Islands} & SAT, SMT, and superposition adapters & Supported translation and proof rules; no raw oracle acceptance\\
\code{Forge.Replay} & Proof minimization and audit & Exact original goal, no remaining metavariables, approved axioms\\
\bottomrule
\end{tabularx}
\end{center}

Reification deserves particular attention. To use a rational polynomial
certificate, store a map $\rho$ from polynomial variables to Lean expressions
and a proof that the original term equals the polynomial's denotation at
$\rho$. For an ordered-field inequality, every coefficient cast and operation
must agree with that denotation. An algebraically correct certificate for the
wrong denotation is not a proof of the original statement.

\subsection{Libraries and algorithms with specific roles}

The initial Lean dependency should be \textbf{core Lean plus mathlib}, with
Aesop used through the version selected by the pinned dependency graph when
available. Reuse \code{ring} or \code{ring\_nf} for polynomial identities,
\code{norm\_num} for rational constants, \code{positivity} for cheap sign facts,
and \code{linarith}/\code{nlinarith} for arithmetic leaves and selected
certificate reconstruction \cite{ring,positivity,linarith}. Do not invoke all
of them blindly on every node.

The research implementation uses Python's \code{fractions.Fraction} and a
small sparse-polynomial representation for checking, SciPy/HiGHS for LP
proposals, and SymPy for exact support solving and interpolation. Exact tested
versions are recorded with the results. NumPy constructs numerical matrices;
its floating-point values never decide certificate acceptance.

Optional integrations are Duper for superposition and Lean-SMT with its
supported external solver for SMT proofs. Keep them out of the minimal build
until compatibility tests pass. A numerical semidefinite solver can later
supply Gram-matrix proposals, but an exact replay path is a prerequisite.
An existing e-graph package is not required: \grind\ already maintains
congruence machinery. In particular, an archived experimental Lean equality
saturation project should not be made a critical dependency simply because
its name sounds aligned with the design.

\subsection{Resource allocation and proof-cost awareness}

After the protected baseline attempt, a practical starting policy allocates
local bounded actions rather than unbounded calls: demand expansion, a small
cone, a shallow induction candidate, a small witness grammar, and an eligible
solver island. Rank actions by observed recent payoff, target relevance, cost
estimate, and age. Reserve a periodic fair slot for postponed applicable
actions. The precise weights are tunable engineering parameters, not
probabilistically calibrated guarantees.

For an implementable non-learned initial scheduler, use an eight-slot round
with three demand slots, two cone slots, one structural slot, one witness slot,
and one eligible external-island slot. Empty slots are donated to the oldest
applicable action. Within a slot, rank target-linked obligations before generic
ones, then prefer lower estimated cost; use a deterministic identifier to break
ties. Every fourth round give the oldest postponed action priority. This is a
proposed starting policy, not a measured optimum or a finite-budget guarantee.

Initial per-action limits can be explicit: expand at most 32 theorem instances
per demand quantum; try cone degrees 2, 4, and 6 with at most 4,000 columns;
try at most two induction candidates per structural quantum and induction depth
at most two; enumerate at most 256 witness proposals per quantum with term size
at most six. Track external jobs separately under wall-clock and output-size
limits. Every quantum also consumes the common parent budget, so nested actions
cannot multiply the allowed work by repeatedly resetting their local limit.

Keep a monotone context epoch alongside each immutable exported problem. An
asynchronous result is reconstructed only after checking that its expression
mapping and hypothesis dependencies still apply. Never reuse an old mutable
metavariable assignment merely because its printed goal matches the current
one. If transport to the new context cannot be justified, discard the result
or rerun the action. These rules are especially important when induction,
backtracking, and external jobs overlap.

Count the cost of \emph{reconstruction}. An apparently fast oracle that emits
an enormous proof can lose to a slower search with a small certificate. Use
support minimization, common subexpression sharing, and proof-DAG trimming.
Keep unsat-core or used-premise information as a proposal for minimization,
then replay the reduced proof to verify that no necessary dependency was
removed. Hash-consing expressions reduces duplication but does not excuse
missing scoping checks.

All external requests require limits on time, memory, serialized input/output
size, integer bit length, and parser depth. Do not evaluate returned text as
code. A well-formed but gigantic certificate is a resource risk even when its
mathematics is harmless. Permit cancellation between checker steps and return
an unresolved result on resource exhaustion.

\subsection{User-facing syntax and explanations}

The following is \textbf{proposed syntax}, not syntax accepted by the delivered
package:
\begin{Code}
by
  forge (mode := conservative)
    (inductionDepth := 2)
    (polyDegree := 4)
    (witnessSize := 6)
    (external := false)
\end{Code}

A companion \code{forge?} should suggest a stable replay script or an explicit
certificate invocation. Users should see which premises were used, which
induction variable was selected, which parameters were generalized, and what
witness was constructed. Failure diagnostics should distinguish a missing
sign guard, a degree cap, unsupported reification, an unproved generalized
lemma, an unavailable backend, and an exhausted search budget.

The report's generated Lean files are a narrower first artifact: explicit
proofs for saved certificates. They contain no search implementation of
\forge. Hiding that distinction behind an umbrella import would misrepresent
the state of the project.

\section{Soundness obligations for the whole system}

\begin{proposition}[Compositional soundness target]
Suppose every fact node contains a Lean proof in its recorded context; every
refinement node is justified by an accepted Lean theorem or recursor; every
certificate adapter proves its exact translated conclusion; and every scope
transfer is checked. Then a closed root of the acyclic proof graph supplies a
proof of the original goal, independently of how search candidates were
selected.
\end{proposition}
\begin{proof}
Proceed in topological order through the proof graph. Base facts are valid by
assumption. At a theorem-application node, apply its proof to the previously
constructed premise proofs. At a refinement node, combine the checked branch
proofs with its theorem or recursor. At a certificate node, use the adapter's
proved conclusion. Scope transfer abstracts and instantiates only validated
dependencies. The final node therefore inhabits the original proposition.
\end{proof}

This is a design-level argument, not a mechanized proof of the present Python
implementation. Lean integration must instantiate each hypothesis of the
argument. An unsound term translator is not repaired by an exact polynomial
checker. A sound checker is not repaired by a replay path that assumes its
result through an oracle axiom. An acyclic certificate DAG is not sufficient
if a proof parent was imported from the wrong branch.

Testing should deliberately attack those boundaries. Mutate a denominator
condition; swap two polynomial variables; replace a strict inequality with a
weak one; change an exponent; reverse an equality transport; remove a branch
from a box cover; alter a typeclass instance; point a proof node to a future
node; and replay under a different local context. The intended result is
rejection, not a best-effort reinterpretation of the malformed evidence.
